\section{Introduction}

En séquentiel, on parle de complexité en \notion{temps} et en \notion{espace}, souvent dans le pire cas. La complexité en temps est le nombre d'opérations élémentaires effectuées par un algorithme, tandis que la complexité en espace est la quantité de mémoire utilisée par l'algorithme.


En algorithmique distribuée, on parle de complexité en \notion{temps} et en \notion{messages}. La complexité en temps est le nombre de tours de communication nécessaires pour que tous les nœuds du réseau terminent l'exécution de l'algorithme. La complexité en messages est le nombre total de messages échangés entre les nœuds pendant l'exécution de l'algorithme. 

\begin{definition}{}{complexité en temps}
  La complexité en temps d'un algorithme distribué est le nombre de tours de communication nécessaires pour que tous les nœuds du réseau terminent l'exécution de l'algorithme, dans le pire cas.
\end{definition}

\begin{definition}{}{temps asynchrone}
  $A$ un aglorithme distribué sur $G$ un graphe d'infrastucture, $\text{TEMPS}(AG)$ est le nombre d'\notion{unités de temps} écoulées jusqu'au dernier message \notion{consommé} dans le pire des cas (sur toutes les entrées et sur \notion{tous les scénarios}). \par
  L'unité de temps étant la durée maximale pour qu'un message traverse une arête.

\end{definition}

\begin{definition}{}{système synchrone, systèe me asynchrone}
  Un système est dit \notion{synchrone} si les nœuds disposent d'une horloge globale et que les messages sont transmis en un temps borné. \par
  Un système est dit \notion{asynchrone} si les nœuds ne disposent pas d'une horloge globale et que les messages peuvent être transmis en un temps non borné.
\end{definition}

\begin{definition}{}{complexité en messages}
$A$ un aglorithme distribué sur $G$ un graphe d'infrastucture, $\text{MESSAGES}(AG)$ est le nombre de messages échangés dans le pire des cas (sur toutes les entrées et sur \notion{tous les scénarios}). \par

\end{definition}

